% Laget automatisk av elevark/lag_elevark.py fra opplegg/transformator.html. Endringer her blir overskrevet.
\documentclass[11pt]{article}
\usepackage{elevark}
\begin{document}
\elevarkhode{Fysikk 2 · Elektromagnetisk induksjon}{Transformatoren}{https://fysikklab.no/opplegg/transformator.html}{Opplegget på nett}

En transformator er to spoler på en felles jernkjerne. Vekselstrøm i den ene spolen gir en fluks som hele tiden endrer seg, og den induserer en spenning i den andre. Med flere vindinger på sekundærspolen blir spenningen høyere. Stemmer forholdet $U_{\mathrm{s}}/U_{\mathrm{p}} = N_{\mathrm{s}}/N_{\mathrm{p}}$?

\section*{Hensikt}

\begin{itemize}
\item Undersøke sammenhengen mellom spenningene og vindingstallene i en transformator.
\item Se at en transformator bare virker med vekselspenning.
\end{itemize}

\section*{Utstyr}

\begin{itemize}
\item Transformatorsett med U-formet jernkjerne og et lukkeåk, og spoler med for eksempel 100, 200, 400, 800 og 1600 vindinger.
\item Spenningskilde som gir lav vekselspenning (rundt 6~V), og to multimetre innstilt på vekselspenning (V\textasciitilde{}).
\item Bruk bare lav spenning. Sekundærspenningen kan bli høyere enn primærspenningen, så ikke ta på kontaktene mens spenningen er på.
\end{itemize}

\section*{Teorien}

Vekselstrømmen i primærspolen gir en fluks $Φ$ i kjernen som endrer seg hele tiden. Den samme fluksen går gjennom begge spolene. I primærspolen er $U_{\mathrm{p}} \approx N_{\mathrm{p}}ΔΦ/Δt$, og i sekundærspolen er $U_{\mathrm{s}} = N_{\mathrm{s}}ΔΦ/Δt$. Det gir $U_{\mathrm{s}}/U_{\mathrm{p}} = N_{\mathrm{s}}/N_{\mathrm{p}}$. I en ideell transformator går ingen energi tapt, så $U_{\mathrm{p}}I_{\mathrm{p}} = U_{\mathrm{s}}I_{\mathrm{s}}$.

\section*{Framgangsmåte}

\begin{enumerate}
\item Sett primærspolen (for eksempel 400 vindinger) på den ene foten av kjernen og en sekundærspole på den andre. Legg på åket.
\item Koble primærspolen til vekselspenningen og mål $U_{\mathrm{p}}$. Mål $U_{\mathrm{s}}$ over sekundærspolen.
\item Bytt sekundærspolen og gjenta for alle vindingstallene. Skriv inn i tabellen.
\item Ta av åket og mål på nytt. Hva skjer?
\item Bytt til likespenning med samme størrelse. Hva viser voltmeteret over sekundærspolen?
\end{enumerate}

\section*{Oppgaver}

\begin{oppgaveliste}
\oppgave \fase{gjett}{Gjett}Primærspolen har 400 vindinger og 6,0~V. Hva tror du sekundærspenningen blir med 1600 vindinger?
\svar
\oppgave \fase{test}{Mål}Gjør målingene, eller trykk på «Fyll inn eksempel». Ligger punktene på en rett linje? Hvor stor del av den ideelle spenningen får dere?
\svar
\oppgave \fase{forklar}{Forklar}Hvorfor blir sekundærspenningen mye mindre når åket tas av? Hvorfor virker ikke transformatoren med likespenning?
\svar
\oppgave \fase{test}{Regn}En lader omformer 230~V til 12~V og gir 36~W. Primærspolen har 1150 vindinger. Hvor mange vindinger har sekundærspolen, og hvor store er strømmene på hver side (ideell transformator)?
\svar
\oppgave \fase{forklar}{Vurder}Hvor blir det av energien som går tapt i en ekte transformator?
\svar
\oppgave \fase{utfordring}{Utfordring}Et kraftverk sender 100~MW gjennom en linje med resistans 5,0~Ω. Regn ut effekttapet i linja når spenningen er 400~kV, og når den er 20~kV. Hva forteller det om hvorfor strømnettet bruker transformatorer?
\svar
\end{oppgaveliste}

\end{document}
